
  \cventry
    {Perk} % Organization
    {Software Engineer} % Job title
    {Edinburgh, Scotland (Hybrid)} % Location
    {Jun 2024 – Present} % Date(s)
    {
      \begin{itemize}
        \item Led delivery of the company-wide Perk Rebrand across the Invoicing \& VAT domain against a fixed, non-negotiable deadline; designed the GroupDomain feature switch enabling independent squad releases, achieving zero critical regressions at launch.
        \item Diagnosed and resolved a production incident where async timing gaps caused 2,970+ customer invoices to incorrectly omit credit line items (a 110$\times$ increase post-deploy); designed an event-driven invoice issuance flow with SQS queue delay, raised Terraform infrastructure PRs, and coordinated cross-team deployment — fully resolving the root cause.
        \item Applied Monte Carlo simulations to reduce AI benchmarking runtime and OpenAI API costs by \textasciitilde60\%; built Datadog dashboards enabling data-led prompt optimisation, improving invoice verification pass rates from 84\% to 98\% on DACH customer data.
        \item Proactively authored a technical design document for Perk's Deferred Payments domain — a technically complex post-stay reconciliation system that determines the correct total charge for hotel stays after unknown extras are confirmed — secured stakeholder review and sign-off without top-down direction, transforming a high-toil legacy system into a funded engineering workstream now in execution.
        \item Delivered Self-Serve Supplier Invoices MVP with strict scope management, embedded feature adoption metrics from launch, and cross-functional alignment across engineering, product, and design.
        \item Designing and delivering service decomposition work as part of Perk's strategic migration from a large-scale Django monolith to microservices; applying distributed system patterns across service boundaries, inter-service API design, and data ownership.
        \item Mentored new engineers and authored the onboarding guide for Perk's Pace Sustainable Engineering (PSE) programme — a proactive system health and risk initiative — adopted into official team documentation.
      \end{itemize}
    }
